% !TEX root = ../hw02.tex
\leavevmode
\begin{enumerate}
  \item \textit{Newton's method.}
    For $x>0$, let
    \[
      F(x)=x-a-\frac1x,\qquad F'(x)=1+\frac1{x^2}>0.
    \]
    Newton's method for $F(x)=0$ is therefore
    \[
      \begin{aligned}
        x_{n+1}
          &=x_n-\frac{F(x_n)}{F'(x_n)}\\
          &=x_n-\frac{x_n-a-1/x_n}{1+1/x_n^2}\\
          &=\frac{x_n(1+x_n^2)-(x_n^3-ax_n^2-x_n)}{1+x_n^2}\\
          &=\frac{ax_n^2+2x_n}{1+x_n^2},
      \end{aligned}
    \]
    which is the stated iteration. If $x_n>0$, then the numerator
    and denominator are positive. Since $x_0>0$, induction shows
    that every iterate is positive and the algorithm is well defined.

  \item \textit{Identification of the limit.}
    Suppose $x_n\to L$. Positivity implies $L\geq0$; we first
    rule out $L=0$. If $L=0$, there is an $N$ such that
    $0<x_n<1$ for all $n\geq N$. For these indices,
    \[
      \frac{x_{n+1}}{x_n}
        =\frac{ax_n+2}{1+x_n^2}>1.
    \]
    The tail is then increasing and satisfies $x_n\geq x_N>0$,
    contradicting $x_n\to0$. Hence $L>0$.

    Passing to the limit in the recurrence gives
    \[
      L=\frac{aL^2+2L}{1+L^2}.
    \]
    Thus $L(L^2-aL-1)=0$. Since $L>0$, we obtain
    \[
      L^2-aL-1=0,\qquad
      L=\frac{a+\sqrt{a^2+4}}2.
    \]
    This is the unique positive solution of $L=a+1/L$.
    The continued fraction $b$ is a positive fixed point of
    $f(x)=a+1/x$, so uniqueness implies $L=b$.

  \item \textit{Bounds when $a=3/2$.}
    For the remaining parts, set $a=3/2$ and write
    \[
      g(x)=\frac{\tfrac32x^2+2x}{1+x^2},\qquad x>0.
    \]
    We have $g(x)>0$ and
    \[
      2-g(x)
        =\frac{2(1+x^2)-\tfrac32x^2-2x}{1+x^2}
        =\frac{(x-2)^2}{2(1+x^2)}\geq0.
    \]
    The nonnegativity of the numerator is illustrated in
    Figure~\ref{fig:hw02:nonnegative-parabola}.
    \begin{marginfigure}
      \centering
      % !TEX root = ../../problems/hw02.tex
% IMG_9094: y = (x - 2)^2, with axes and intercepts made explicit.
\begin{tikzpicture}[x=0.6cm,y=0.38cm,>=Stealth,font=\footnotesize]
  \draw[<->,black!65] (-1.05,0) -- (5,0) node[right] {$x$};
  \draw[<->,black!65] (0,-0.85) -- (0,7) node[above] {$y$};
  \draw[<->,line width=0.65pt,domain=-0.55:4.55,samples=100,smooth]
    plot (\x,{(\x-2)^2});
  \fill (0,4) circle[radius=1.6pt];
  \fill (2,0) circle[radius=1.6pt];
  \node[below left,inner sep=2pt] at (0,0) {$0$};
  \node[left=3pt] at (0,4) {$4$};
  \node[below=3pt] at (2,0) {$2$};
\end{tikzpicture}

      \caption{The numerator $(x-2)^2$ is nonnegative,
        with equality only at $x=2$.}
      \label{fig:hw02:nonnegative-parabola}
    \end{marginfigure}
    Therefore $0<g(x)\leq2$ for every $x>0$, with $g(x)=2$
    if and only if $x=2$. Applying this to $x_n>0$ proves
    \[
      0<x_n\leq2\qquad\text{for every }n\geq1,
    \]
    regardless of whether $x_0\leq2$.

    % Keep the monotonicity proof and its margin plot on the same page.
    \newpage

  \item \textit{Monotonicity.}
    If $x_0=2$, then $g(2)=2$, so induction gives $x_n=2$
    for all $n\geq0$.

    Now suppose $x_0\ne2$. The equality condition in part (c)
    gives $0<x_1<2$. Repeatedly applying the same condition yields
    $0<x_n<2$ for every $n\geq1$. For these indices,
    \[
      \begin{aligned}
        x_{n+1}-x_n
          &=\frac{\tfrac32x_n^2+2x_n-x_n(1+x_n^2)}{1+x_n^2}\\
          &=\frac{x_n(-2x_n^2+3x_n+2)}{2(1+x_n^2)}\\
          &=\frac{x_n(2-x_n)(2x_n+1)}{2(1+x_n^2)}>0.
      \end{aligned}
    \]
    Every factor in the last expression is positive. Equivalently,
    the quadratic $-2x^2+3x+2=(2-x)(2x+1)$ is positive on
    $(-1/2,2)$, as shown in Figure~\ref{fig:hw02:concave-parabola}.
    \begin{marginfigure}
      \centering
      % !TEX root = ../../problems/hw02.tex
% IMG_9096: y = -2x^2 + 3x + 2, with axes and intercepts made explicit.
\begin{tikzpicture}[x=0.85cm,y=0.4cm,>=Stealth,font=\footnotesize]
  \draw[<->,black!65] (-1.3,0) -- (2.95,0) node[right] {$x$};
  \draw[<->,black!65] (0,-3.1) -- (0,4.2) node[above] {$y$};
  \draw[<->,line width=0.65pt,domain=-0.95:2.45,samples=100,smooth]
    plot (\x,{-2*(\x)*(\x)+3*(\x)+2});
  \foreach \xx/\yy in {-0.5/0,2/0,0/2} {
    \fill (\xx,\yy) circle[radius=1.6pt];
  }
  \node[below right,inner sep=2pt] at (0,0) {$0$};
  \node[above left=3pt] at (-0.5,0) {$-\tfrac12$};
  \node[above right=3pt] at (2,0) {$2$};
  \node[left=3pt] at (0,2) {$2$};
\end{tikzpicture}

      \caption{The quadratic $-2x^2+3x+2$ is positive between
        its roots $-1/2$ and $2$.}
      \label{fig:hw02:concave-parabola}
    \end{marginfigure}
    Thus $(x_n)_{n\geq1}$ is strictly increasing when $x_0\ne2$.
    The restriction $n\geq1$ is necessary because if $x_0>2$, the first
    iterate is below $2$ and hence below $x_0$.

  \item \textit{Global convergence.}
    If $x_0=2$, the sequence is constant and converges to $2$.
    If $x_0>0$ and $x_0\ne2$, parts (c) and (d) show that
    $(x_n)_{n\geq1}$ is increasing and bounded above by $2$.
    The monotone convergence theorem therefore gives a limit.
    By part (b), with $a=3/2$, this limit is
    \[
      \frac{\tfrac32+\sqrt{(\tfrac32)^2+4}}2
        =\frac{\tfrac32+\tfrac52}2=2.
    \]
    Adding the initial term does not affect convergence, so
    $x_n\to2$ for every allowed initial state $x_0>0$.
\end{enumerate}
